\documentclass[12pt]{article}
\usepackage{amsfonts,amsthm,amsmath,amssymb,mathtools}
\usepackage{mathrsfs}
\usepackage{enumitem}
\usepackage{tikz-cd}
\usepackage{geometry}
\usepackage{mlmodern}

\geometry{letterpaper, portrait, margin=1in}

\setcounter{section}{0}

\renewcommand{\thesection}{\S\arabic{section}}
\renewcommand{\thesubsection}{\arabic{section}}
\newcommand{\x}{\mathbf{x}}

\newtheorem{thm}{Theorem}
\newenvironment{theorem}[1]{
	\renewcommand{\thethm}{#1}
	\begin{thm}
		}{\end{thm}}

\newtheorem{lma}{Lemma}
\newenvironment{lemma}[1]{
	\renewcommand{\thelma}{#1}
	\begin{lma}
		}{\end{lma}}

\theoremstyle{definition}
\newtheorem{ex}{}
\newenvironment{exercise}[1]{
	\renewcommand{\theex}{#1}
	\begin{ex}
		}{\end{ex}}

\title{Homework 4}
\author{Connor Mooneyhan}
\date{February 12, 2026}

\begin{document}

\maketitle

\begin{ex}
	Recall that a metric space $(X, d)$ is called complete if every Cauchy sequence in $X$ converges.
	Show that if $X$ is compact then it must be complete.
\end{ex}
\begin{proof}
	Let $(a_n)$ be a Cauchy sequence in $X$.
	Since $X$ is a compact metric space, it is sequentially compact and thus $(a_n)$ has a subsequence $(b_n)$ which converges to a point $p \in X$.
	Since $(a_n)$ is Cauchy, it converges to the same $p \in X$, and thus $X$ is complete.

	% Let $\varepsilon > 0$.
	% Then $\exists N$ such that for all $n, m > N$, $d(a_n, a_m) < \varepsilon / 2$, and $\exists M > 0$ such that for all $m > M$, $d(a_{i_m}, p) < \varepsilon / 2$.
	% Now take $L = \max(M, N)$.
	% Then for all $n, m > L$, we have $d(a_{i_n}, p) + d(a_{i_n}, a_m) \leq d(a_m <)$
\end{proof}

\begin{ex}
	Consider $C([0, 1])$ with the $C^0$ metric.
	Let \begin{equation*}
		B = \left\lbrace f \in C([0, 1]) : d(f, 0) \leq 1 \right\rbrace
	\end{equation*}
	where $\mathbf{0}$ is the constant function $g \equiv 0$.
	Note that $B$ is closed and bounded.
	Show that it is not compact.
\end{ex}
\begin{proof}
	Define \begin{equation*}
		f_n(x) = \begin{cases}
			2nx     & \text{if } x \in [0, 1/2n)   \\
			2 - 2nx & \text{if } x \in [1/2n, 1/n) \\
			0       & \text{otherwise.}
		\end{cases}
	\end{equation*}
	Then each $f_n(x) \in C([0,1])$, and the uniform limit of the sequence as $n \to \infty$ is \begin{equation*}
		f(x) = \begin{cases}
			1 & x = 0             \\
			0 & \text{otherwise.}
		\end{cases}
	\end{equation*}
	Thus any subsequence must also converge to $f$, but of course $f$ is not continuous on $[0, 1]$ since it has a jump discontinuity at $x = 0$, so $B$ is not sequentially compact.
\end{proof}

\begin{ex}
	Consider $C([0, 1])$ with the $L^1$ metric.
	Let \begin{equation*}
		B = \left\lbrace f \in C([0, 1]) : d(f, \mathbf{0}) \leq 1 \right\rbrace
	\end{equation*}
	and show that it is not compact.
\end{ex}
\begin{proof}
	Define \begin{equation*}
		f_n(x) = \begin{cases}
			4n^2x     & \text{if } x \in [0, 1/2n)   \\
			4n - 4n^2x & \text{if } x \in [1/2n, 1/n) \\
			0       & \text{otherwise.}
		\end{cases}
	\end{equation*}
  Then each $f_n \in C([0,1])$, and the integral of each $f_n$ over $[0, 1]$ is 1, so each $f_n \in B$.
  However, since these do not converge, there is no convergent subsequence, and thus $X$ is not compact.
\end{proof}

\begin{ex}
	Let $K_1 \supset K_2 \supset K_3 \supset \cdots$ be a sequence of non-empty compact subsets of a metric space $(X, d)$.
	Show that $\bigcap_{n=1}^\infty K_n$ is also non-empty.
\end{ex}
\begin{proof}
  For each $n > 0$, let $a_n \in K_n$.
  Then $(a_n) \subset K_1$, so $(a_n)$ has a subsequence $(a_{i_m})$ which converges to a point $p \in K_1$ by sequential compactness of $K_1$.
  However, for each $n > 0$, $(a_m)_{m = n}^\infty \subset K_n$, so indeed all but finitely many entries of $(a_{i_m})$ also lie in $K_n$, and thus their limit (which is unique amongst the subsequential limits of $(a_{i_m})$) is in $K_n$ by compactness of $K_n$.
  Since this is true for all $n$, we get that $p \in K_n$ for all $n$, and thus $p \in \bigcap K_n$.
\end{proof}

\end{document}
